% =============================================================================
% Project : THE_PHILOSOPHY_OF_VIABILITY
% Book    : Book I — Ontology of Viability
% Book ID : VS.B01
%
% File    : 4_ASSETS/illustrations/
%           VS.B01.P01.C02.S08.SS06.P003.IL01.continuity_through_transformation.tex
%
% Asset   : IL01 — Continuity Through Transformation
% Location: Chapter 02 / Section 08 / Subsection 06 / Paragraph 003
%
% Purpose :
%   To provide a concrete visual illustration of a process
%   whose visible configurations change across successive
%   phases while a proposed connecting relation remains
%   available for examination.
%
% Conceptual Caution :
%   The illustration does not establish numerical identity,
%   material identity, or a universal criterion of continuity.
%
% Requires (preamble):
%   \usepackage{tikz}
%   \usetikzlibrary{arrows.meta}
% =============================================================================

\begin{figure}[htbp]
    \centering

    % horizontal distance between phase centres (panel width 3.4 cm + gap 1.0 cm)
    \def\phpitch{4.4}

    \begin{tikzpicture}[
        % --- panel ---
        phasebox/.style={
            draw=black!60, fill=gray!6, rounded corners=4pt,
            minimum width=3.4cm, minimum height=3.0cm, line width=0.7pt
        },
        % --- drawing elements ---
        flow/.style={line width=2.2pt, line cap=round, draw=blue!55!black},
        bank/.style={line width=0.5pt, draw=black!70},
        % --- arrows and relation line ---
        transition/.style={-{Stealth[length=3mm]}, line width=1.1pt},
        rel/.style={dashed, thick, draw=black!70},
        % --- one phase panel: #1 title / #2 flow rise / #3 bank tilt / #4 label ---
        pics/phase/.style args={#1/#2/#3/#4}{
            code={
                \node[phasebox] (-box) at (0,0) {};
                \node[font=\small\bfseries] at (0,1.15) {#1};
                \draw[flow] (-1.2,-0.6)
                    .. controls (-0.6,{-0.6+#2}) and (0.6,{-0.6+#2})
                    .. (1.2,-0.6);
                \draw[bank] (-1.35,-0.95) -- ({-1.35+#3},-0.35);
                \draw[bank] ( 1.35,-0.95) -- ({ 1.35-#3},-0.35);
                \node[font=\scriptsize] at (0,-1.2) {#4};
            }
        }
    ]

    % ---------------------------------------------------------------
    % Phases
    % ---------------------------------------------------------------
    \pic (A) at (0,0)              {phase=Phase A/0.40/0.00/Initial configuration};
    \pic (B) at (\phpitch,0)       {phase=Phase B/0.90/0.30/Modified configuration};
    \pic (C) at ({2*\phpitch},0)   {phase=Phase C/1.25/0.60/Further configuration};

    % ---------------------------------------------------------------
    % Transitions
    % ---------------------------------------------------------------
    \draw[transition] (A-box.east) -- (B-box.west);
    \draw[transition] (B-box.east) -- (C-box.west);

    % ---------------------------------------------------------------
    % Proposed relation (undirected: no arrowheads)
    % ---------------------------------------------------------------
    \draw[rel] (A-box.south) -- ++(0,-0.6) coordinate (ra);
    \draw[rel] (B-box.south) -- ++(0,-0.6) coordinate (rb);
    \draw[rel] (C-box.south) -- ++(0,-0.6) coordinate (rc);
    \draw[rel] (ra) -- (rc);

    \node[below=4pt of rb, font=\small\itshape] {
        Proposed relation connecting successive phases
    };

    \end{tikzpicture}

    \caption{An illustrative process across changing configurations.
    The visible determinations differ between phases, while a
    connecting relation is proposed for examination. The image
    illustrates a question of continuity; it does not establish
    numerical identity or the preservation of an invariant.}
    \label{fig:continuity-through-transformation}
\end{figure}

% =============================================================================
% END OF ILLUSTRATION IL01
% =============================================================================